\documentclass[11pt,a4paper]{article}
\usepackage[T1]{fontenc}
\usepackage[utf8]{inputenc}
\usepackage{lmodern}
\usepackage{microtype}
\usepackage[margin=27mm,headheight=15pt]{geometry}
\usepackage{amsmath,amssymb,amsthm,mathtools}
\usepackage{booktabs,longtable,array}
\usepackage{xcolor}
\usepackage{listings}
\usepackage{enumitem}
\usepackage{fancyhdr}
\usepackage{xurl}
\usepackage{hyperref}
\definecolor{ink}{RGB}{29,52,75}
\definecolor{muted}{RGB}{78,89,101}
\definecolor{codebg}{RGB}{246,247,249}
\hypersetup{colorlinks=true,linkcolor=ink,citecolor=ink,urlcolor=ink,
  pdftitle={SystematicRadicals: From algebraic numbers to certified radical towers},
  pdfauthor={Prepared for Vladimir Reshetnikov}}
\pagestyle{fancy}
\fancyhf{}
\fancyhead[L]{\small\sffamily\color{muted}SYSTEMATIC RADICALS}
\fancyhead[R]{\small\sffamily\color{muted}Theory and reference implementation}
\fancyfoot[C]{\small\thepage}
\renewcommand{\headrulewidth}{0.3pt}
\setlength{\parskip}{4pt}
\setlength{\parindent}{1.1em}
\setcounter{tocdepth}{2}
\setlist[enumerate]{itemsep=3pt,topsep=5pt}
\lstset{basicstyle=\small\ttfamily,backgroundcolor=\color{codebg},
  frame=single,rulecolor=\color{black!12},breaklines=true,
  columns=fullflexible,keepspaces=true,showstringspaces=false,
  xleftmargin=5pt,xrightmargin=5pt,aboveskip=9pt,belowskip=8pt}
\newtheorem{theorem}{Theorem}[section]
\newtheorem{proposition}[theorem]{Proposition}
\newtheorem{lemma}[theorem]{Lemma}
\newtheorem{corollary}[theorem]{Corollary}
\theoremstyle{definition}
\newtheorem{definition}[theorem]{Definition}
\theoremstyle{remark}
\newtheorem{remark}[theorem]{Remark}
\newcommand{\Q}{\mathbb Q}
\newcommand{\C}{\mathbb C}
\newcommand{\Gal}{\operatorname{Gal}}
\newcommand{\Tr}{\operatorname{Tr}}
\newcommand{\id}{\operatorname{id}}
\newcommand{\rad}{\operatorname{rad}}
\newcommand{\code}[1]{\texttt{#1}}
\newcommand{\status}[1]{\textsf{#1}}
\newcommand{\z}{\zeta}

\begin{document}
\begin{titlepage}
\vspace*{17mm}
{\Large\sffamily\color{ink}SYSTEMATIC RADICALS\par}
\vspace{7mm}
{\huge\bfseries From algebraic numbers to\\[3pt]certified radical towers\par}
\vspace{7mm}
{\large A constructive solver for exact polynomial \code{Root} objects\par}
\vspace{12mm}
{\large Prepared for Vladimir Reshetnikov\par}
{\normalsize Version 1.0.0\quad\textbullet\quad 8 September 2026\par}
\vspace{13mm}
\begin{abstract}
An exact algebraic number need not become a radical expression when passed to a
computer algebra system's usual radical converter. This is an algorithmic
limitation, not a different notion of algebraic solvability. We develop a
systematic replacement consisting of inexpensive structural reductions and a
fully specified general construction. The latter builds an abstract splitting
field, computes its automorphisms, passes to a cyclotomic base, and converts a
prime-factor composition series into a radical tower by exact Lagrange
resolvents. A separate checker validates positive certificates using rational
linear algebra and polynomial identities, without reconstructing the Galois
group. A coherent-branch theorem explains why principal radicals can represent
the entire root set without combinatorial branch guessing.

The accompanying software contains a Wolfram Language interface with native
fast paths and a self-contained Python/SymPy implementation of the general
construction. Both examples in the motivating question have particularly
small formulas: a reciprocal sextic reduces to a cubic, and a solvable quintic
is a Dickson polynomial. Completeness is an unbounded algorithmic statement;
default field-degree and time budgets deliberately allow an inconclusive
resource-limit result.
\end{abstract}
\vfill
\noindent\begin{minipage}{\textwidth}
\small\textbf{Validation status.} The Python implementation passed all 36 supplied
regression tests under Python 3.13.5 and SymPy 1.14.0. They include exact tower
certificates, selected-root checks, a nonsolvable polynomial, and negative
protocol tests. Wolfram Language tests are supplied but were \emph{not executed}:
the available kernel connection failed. No claim of formal proof-assistant
verification is made.
\end{minipage}
\end{titlepage}

\begingroup
\small
\setlength{\parskip}{0pt}
\tableofcontents
\endgroup
\newpage

\section{Problem, scope, and the meaning of success}
\label{sec:scope}
The motivating question asks why \code{ToRadicals} sometimes leaves behind a
\code{Root} although a radical expression exists \cite{question}. The official
documentation explicitly acknowledges such cases and guarantees conversion
through degree four \cite{toradicals}. The supplied notebook
\code{FindExtension.nb} explores ways to discover suitable auxiliary fields.
This article turns that exploratory idea into a specified construction, with
explicit success, impossibility, and resource-limit outcomes.

Our basic input is a fixed number $\alpha\in\overline{\Q}\subset\C$.
Its defining data include an exact polynomial and an embedding, such as a
\code{Root} object. We exclude parameters, approximate coefficients, and
transcendental constants whose algebraicity has not been established. The
Wolfram interface first computes the minimal polynomial over $\Q$, so it also
accepts exact algebraic expressions beyond a syntactically isolated
\code{Root}. The Python polynomial entry point instead requires an irreducible
polynomial over $\Q$; its selected-root entry point performs the factor
selection itself.

\begin{definition}
A \emph{radical expression over $\Q$} is a finite expression built from rational
numbers by addition, subtraction, multiplication, division by nonzero
quantities, and extraction of positive-integer roots. Complex intermediate
values are allowed. Equivalently, we allow rational powers together with the
field operations and $i=\sqrt{-1}$.
\end{definition}
The implementation checks this grammar. A trigonometric value such as
$2\cos(2\pi/7)$ is not accepted as the final answer merely because it is
algebraic. Roots of unity can, however, be written within the grammar:
\begin{equation}\label{eq:zeta-principal}
 \z_m=(-1)^{2/m}=\exp(2\pi i/m),
\end{equation}
where the power on the left has the principal complex meaning. The exponential
on the right is explanatory notation and does not occur in the exported
expression.

It is important to solve the minimal polynomial of the selected number, not
an arbitrary polynomial that happens to vanish there. For example,
\[
 (x-10)(x^5-x-1)
\]
has the easily expressible root $10$ and other roots that are not expressible
by radicals. Testing the whole product for solvability would incorrectly
reject the easy root. Conversely, an irreducible polynomial cannot have one
radical-solvable conjugate and another nonsolvable conjugate.

\begin{theorem}[The criterion for a selected algebraic number]
\label{thm:criterion}
Let $f$ be the minimal polynomial of $\alpha$ over $\Q$ and let $K$ be its
splitting field. Then $\alpha$ has a radical expression over $\Q$ if and only if
$\Gal(K/\Q)$ is solvable.
\end{theorem}
\begin{proof}
The usual Galois solvability theorem states that a polynomial in characteristic
zero is solvable by radicals exactly when its Galois group is solvable
\cite[Theorem 5.34]{milne}. To apply the criterion to one root, observe that if
$\alpha$ belongs to a radical tower, adjoining the necessary roots of unity and
all conjugates of the tower produces a finite normal extension with solvable
Galois group. Every conjugate of $\alpha$ lies in that extension. Its splitting
field therefore has a solvable group, being a quotient of a solvable group.
The converse follows by constructing radicals for all roots. Sections
\ref{sec:closure}--\ref{sec:branches} give that construction explicitly.
\end{proof}

Three distinct questions must not be confused. Is the number solvable by
radicals? Can the present run find a representation within its budget? Is a
particular proposed representation the \emph{selected} root rather than a
different conjugate? The package keeps these questions separate. A timeout is
not a negative answer to the first question, and the identity $f(e)=0$ is not
by itself an answer to the third.

\section{What the experimental notebook contributes}
\label{sec:notebook}
The archive supplied with the request contains one notebook, marked as created
by Wolfram 15.0. We inspected its core definitions, including
\code{findExtension}, \code{findExtension3}, \code{project}, \code{factor},
\code{solve}, \code{fuzz}, and the simplification helpers. The notebook is an
experimental working document with many subsequent examples and saved outputs;
the discussion here concerns those core algorithms, not a claim to have
re-executed every notebook cell.

The extension searches eliminate coefficients from a proposed quadratic or
cubic factor and look for a manageable coefficient field. This is a useful
resolvent idea: coefficients of factors are symmetric functions of subsets of
roots, and their stabilizers identify intermediate fields of a splitting
field. The projection method examines real and imaginary parts of conjugates.
The fuzzing method samples sums of conjugates, approximates them at high
precision, and uses \code{RootApproximant} to seek a lower-degree description.
These are plausible discovery techniques. They do not, by themselves,
constitute a complete decision or construction procedure.

There are several concrete reasons to change the architecture. An unbounded
random search has no negative stopping criterion. Failure to locate a
quadratic or cubic auxiliary field says nothing about a needed cyclic
extension of larger prime degree. A candidate factor must be checked by exact
polynomial division and a selected root by exact algebraic equality; the
notebook's uses of \code{PossibleZeroQ} and numerical recognition are not the
acceptance condition used here. In addition, two early definitions refer to
\code{Simplify`SimpifyCount}, whereas later definitions use
\code{Simplify`SimplifyCount}. This spelling inconsistency is not a dependable
complexity metric unless the first symbol was separately defined. Depending on
an undocumented internal simplification symbol would be fragile even with
consistent spelling.

The new implementation does not perpetuate a randomized search under a larger
time limit. It replaces it by exact structural identities and then by a
finite group-and-field construction. Simplicity of presentation is a separate
optimization concern. A simple expression is preferred when a fast identity
provides one; otherwise a correct straight-line radical program is more
important than a cosmetically attractive but unverified formula.

A further distinction is between subfields of $\Q(\alpha)$ and subfields of its
normal closure. A prime-degree root field has no proper intermediate fields,
yet its normal closure can contain precisely the auxiliary fields needed to
solve it. Restricting extension discovery to the root field is therefore not
complete. The normal-closure construction below avoids that restriction.

\section{Structural methods before Galois computation}
\label{sec:fast}
For many examples the best solution is to avoid constructing a large splitting
field. Both implementations recognize exact identities. Recognition uses
rational coefficient comparisons, not numerical fitting.

\subsection{Translation and power composition}
A monic polynomial
\[
 f(x)=x^n+a_{n-1}x^{n-1}+\cdots
\]
can first be centered by $x=y+c$, $c=-a_{n-1}/n$. If all nonconstant exponents
of the centered polynomial have common divisor $m>1$, then
\[
 f(y+c)=g(y^m).
\]
For every root $b$ of $g$, the corresponding roots are
\begin{equation}\label{eq:power-lift}
 c+\z_m^k b^{1/m},\qquad 0\leq k<m.
\end{equation}
This includes shifted binomials and compositions of several such reductions.
The equation is exact for every listed candidate. When the original input is
irreducible in characteristic zero its roots are distinct. Repeated roots in
an intermediate reduction do not create a correctness problem, although the
minimal-polynomial preprocessing removes that issue for the top-level input.

Degree at most four is handled by the ordinary exact quadratic, cubic, and
quartic routines. Their results still pass the radical-grammar check. The
structural layer has a finite recursion bound; a miss simply transfers control
to the general construction. The bound is not part of the completeness claim
for the general layer.

\subsection{Dickson polynomials}
Define
\[
 D_0(t,a)=2,\quad D_1(t,a)=t,\quad
 D_n(t,a)=tD_{n-1}(t,a)-aD_{n-2}(t,a).
\]
Induction on $n$ gives
\begin{equation}\label{eq:dickson}
 D_n(u+a/u,a)=u^n+(a/u)^n \qquad (u\ne0).
\end{equation}
Suppose the centered polynomial is $D_n(t,a)+b$. Set
\begin{equation}\label{eq:dickson-radicals}
 v=\frac{-b+\sqrt{b^2-4a^n}}2,\qquad u=v^{1/n}.
\end{equation}
If the displayed $v$ is zero, use the other quadratic root. The candidates are
\begin{equation}\label{eq:dickson-roots}
 c+\z_n^k u+\frac{a}{\z_n^k u},\qquad 0\leq k<n.
\end{equation}
Indeed, $v+a^n/v=-b$, and \eqref{eq:dickson} applies. For a separable polynomial
these give the full root set: a collision between different indices forces the
degenerate situation $v^2=a^n$, or equivalently the vanishing of the quadratic
discriminant. The degenerate cases are excluded for the irreducible
higher-degree inputs to which this path applies. The case $a=0$ is already a
power-composition case.

The parameter $a$ is not searched. For $n\geq3$ the coefficient of $t^{n-2}$ in
$D_n(t,a)$ is $-na$, so it determines $a$ immediately. The full coefficient
identity is then checked. This is a narrow but genuinely arbitrary-degree
family of radical-solvable polynomials.

\subsection{Weighted reciprocal reduction}
Let a monic degree-$2m$ polynomial have coefficients
\[
 f(x)=\sum_{j=0}^{2m} c_jx^j.
\]
Suppose a nonzero rational $a$ satisfies
\begin{equation}\label{eq:weighted}
 c_{m-k}=a^k c_{m+k}\qquad(1\leq k\leq m).
\end{equation}
The recurrence above also gives
\[
 x^k+a^kx^{-k}=D_k(x+a/x,a).
\]
Consequently
\begin{equation}\label{eq:reciprocal}
 f(x)=x^m g(x+a/x),\qquad
 g(t)=c_m+\sum_{k=1}^m c_{m+k}D_k(t,a).
\end{equation}
After solving $g(t)=0$, solve $x^2-tx+a=0$:
\begin{equation}\label{eq:quadratic-lift}
 x=\frac{t\pm\sqrt{t^2-4a}}2.
\end{equation}
No root at zero has been lost, since the constant coefficient in
\eqref{eq:weighted} is nonzero.

The implementation finds candidate $a$ from the first nonzero ratio
$c_{m-k}/c_{m+k}=a^k$, tests whether that ratio is an exact rational $k$th power,
and verifies all of \eqref{eq:weighted}. For even $k$ it considers both rational
signs. It tries both the original and centered polynomial. This is not a
complete search over arbitrary rational substitutions; it is an inexpensive,
proved recognition method with a precise domain.

\section{The two original examples, explicitly}
\label{sec:examples}
\subsection{The sextic is a cubic followed by a quadratic}
For
\[
 f(x)=x^6+x^4-x^3-x^2-1
\]
there is the identity
\begin{equation}\label{eq:original-sextic}
 f(x)=x^3\left((x-x^{-1})^3+4(x-x^{-1})-1\right).
\end{equation}
Thus $t=x-x^{-1}$ satisfies $t^3+4t-1=0$. Define
\begin{equation}\label{eq:t-sextic}
 t=\sqrt[3]{\frac{9+\sqrt{849}}{18}}
   -\sqrt[3]{\frac{\sqrt{849}-9}{18}}.
\end{equation}
Both cube-root radicands are positive, so there is no principal-versus-real
cube-root ambiguity in this formula. If they are denoted $u^3$ and $v^3$, then
$u^3-v^3=1$ and $uv=4/3$. Hence $(u-v)^3+4(u-v)=1$.
The derivative $3t^2+4$ is positive on the real line, so this is the unique real
root of the cubic. The two real roots of the sextic are therefore
\begin{equation}\label{eq:alpha-sextic}
 \alpha_\pm=\frac{t\pm\sqrt{t^2+4}}2.
\end{equation}
Their signs are respectively positive and negative. In particular,
\[
 \boxed{\operatorname{Root}[\#^6+\#^4-\#^3-\#^2-1\,\&,2]
       =\frac{t+\sqrt{t^2+4}}2}
\]
with $t$ from \eqref{eq:t-sextic}; numerically this is about $1.1306854$.
The other four roots follow by inserting the other two cubic roots into
\eqref{eq:alpha-sextic}. The software uses the general reciprocal identity,
not a hard-coded entry for this polynomial.

\subsection{The quintic is a Dickson polynomial}
The second example is
\[
 5x^5-25x^3+25x+6=5\left(D_5(x,1)+\frac65\right).
\]
Let
\[
 u=\left(\frac{-3+4i}{5}\right)^{1/5},\qquad
 \z_5=(-1)^{2/5}.
\]
The five roots are
\begin{equation}\label{eq:quintic-roots}
 \beta_k=\z_5^k u+(\z_5^k u)^{-1},\qquad 0\leq k<5.
\end{equation}
Since $|(-3+4i)/5|=1$, these roots are real. The argument of the principal $u$
lies between $0$ and $\pi/5$, making $\beta_0$ the largest root. Thus the
question's fifth real root can be written particularly simply as
\begin{equation}\label{eq:beta-simple}
 \boxed{\beta_0=
 \left(\frac{-3+4i}{5}\right)^{1/5}+
 \left(\frac{-3-4i}{5}\right)^{1/5}}.
\end{equation}
Its approximate value is $1.80706$. The occurrence of complex intermediates is
allowed by our definition, even though all five outputs are real.

These derivations establish polynomial membership analytically. The selected
root is then identified by its real ordering, and the supplied Python tests
also check the correspondence using the library's root-separation comparison.
Numerical values are for orientation, not the acceptance criterion.

\section{Constructing the splitting field without numerical embeddings}
\label{sec:closure}
Fast paths cannot cover every solvable Galois group. We now describe the
implemented general fallback. The starting point is an irreducible monic
$f\in\Q[x]$, of degree $n\geq2$.

\subsection{Abstract fields and exact factorization}
A field is represented as
\[
 F=\Q[t]/(P(t)),
\]
where $P$ is irreducible. An element is a rational coefficient vector in the
basis $1,t,\ldots,t^{d-1}$, $d=\deg P$. We do not identify $t$ with an
approximate complex root. Addition, multiplication, inversion, equality, and
polynomial evaluation take place in this quotient field.

Start with $P=f$ and factor $f$ over $F$. If all factors are linear, the field
already contains every root and is a splitting field. Otherwise choose a
nonlinear irreducible factor $h(y)\in F[y]$, adjoin one root, and repeat. Exact
number-field factorization is delegated to SymPy's algebraic-field polynomial
arithmetic; the primitive-element change of coordinates is implemented here.
The documented abstract algebraic-field interface allows a defining polynomial
with a formal root symbol \cite{sympyfields}. This avoids trying to build a
primitive element by repeatedly comparing high-precision complex roots.

\subsection{A finite primitive-element search by linear algebra}
Suppose $F=\Q[t]/P$ has degree $d$, and $h(y)$ is monic irreducible over $F$ of
degree $m>1$. The quotient $A=F[y]/h$ is a field of degree $N=dm$ over $\Q$.
Use the basis
\[
 t^iy^j\quad(0\leq i<d,\ 0\leq j<m).
\]
For $c=1,2,3,\ldots$, put $w=t+cy$ and form the rational matrix
\begin{equation}\label{eq:primitive-matrix}
 B_c=\bigl[\operatorname{vec}(1)\ \operatorname{vec}(w)\ \cdots\
             \operatorname{vec}(w^{N-1})\bigr].
\end{equation}
If $B_c$ is nonsingular, solve for the coordinates of $w^N$ and $t$ in this
basis of powers. Writing $w^N=\sum_{j=0}^{N-1}b_jw^j$ gives the new defining
polynomial
\[
 P_{\rm new}(W)=W^N-\sum_{j=0}^{N-1}b_jW^j.
\]
The coordinate expression $t=T(w)$ also gives $y=(w-T(w))/c$. The implementation
checks the old defining relation and the adjoined-root relation after this
substitution.

\begin{lemma}\label{lem:primitive}
The search for a nonsingular $B_c$ terminates, and the resulting polynomial is
the irreducible minimal polynomial of a primitive element of $A$.
\end{lemma}
\begin{proof}
Since the characteristic is zero, $A$ has $N$ distinct embeddings into an
algebraic closure. For any two different embeddings $\sigma,\tau$, the equality
\[
 \sigma(t)+c\sigma(y)=\tau(t)+c\tau(y)
\]
holds for at most one $c$, unless both images of $t$ and $y$ coincide. The
latter would make the embeddings identical, because $t,y$ generate $A$.
Only finitely many $c$ are therefore bad. For every other $c$, $w$ has $N$
distinct conjugates and degree $N$. Its first $N$ powers form a basis, so
$B_c$ is invertible. Conversely, invertibility itself certifies that $w$ has
degree $N$. Its displayed degree-$N$ relation must be its minimal polynomial.
\end{proof}

\subsection{Why adjoining roots stops}
Every intermediate field embeds into a splitting field of $f$: it is generated
by some roots of $f$, and adjoining a root of an irreducible factor over that
field amounts to adjoining another root of $f$. Each proper adjunction
increases the degree. The final degree is at most $n!$, and after finitely many
adjunctions $f$ splits. A splitting field of a separable polynomial is normal.
This gives a finite construction, even though the resulting degree and
coefficient heights may be very large.

The code chooses a factor of smallest available nonlinear degree, and uses a
deterministic tie-breaker. This is a useful local heuristic, not a proof of
minimum intermediate coefficient height. No particular complex ordering of
roots is involved at this stage.

\section{Cyclotomic base change and an explicit Galois group}
\label{sec:group}
Let $K$ be the splitting field just constructed, and write $D=[K:\Q]$. Define
\[
 M=\rad(D)=\prod_{p\mid D}p,\qquad C=\Q(\z_M),\qquad L=KC.
\]
To construct $L$, factor the cyclotomic polynomial $\Phi_M$ over $K$. If it has
a linear factor, a primitive $M$th root already belongs to $K$. Otherwise
adjoin a root of an irreducible factor by the same quotient-field procedure.
The resulting field contains all powers of that root, so $\Phi_M$ splits
there. Both $K$ and $C$ are normal over $\Q$; hence $L$ is normal too.

Only the squarefree product of the prime divisors of $D$ is required, rather
than a root of unity of order $D!$. The reason is that we will use a composition
series with \emph{prime} cyclic factors. This can still be an expensive base
change: the degree of $L$ is bounded by $D\varphi(M)$, not by $D$ alone.

\subsection{Recovering automorphisms from roots of the defining polynomial}
Write $L=\Q[t]/P$ and $d=[L:\Q]$. Factor $P$ over $L$. Normality means that its
$d$ roots $a_1,\ldots,a_d$ all lie in $L$. They give all automorphisms
\[
 \sigma_j(t)=a_j,\qquad
 \sigma_j\!\left(\sum_kq_kt^k\right)=\sum_kq_ka_j^k.
\]
The implementation retains the automorphisms fixing the chosen primitive root
$z$ of $\Phi_M$. They form
\[
 H=\Gal(L/C),\qquad |H|=d/\varphi(M).
\]
It checks the expected cardinality and builds their multiplication table by
exact substitution. A regular permutation representation of that table is then
passed to the finite-group routines for solvability and composition series
\cite{sympygroups}. In particular, this is not a call to an unimplemented
``general Galois oracle.'' The group elements are actually constructed as
field automorphisms.

For degrees through six a preliminary call to SymPy's public
\code{galois\_group} can reject nonsolvable inputs cheaply. The documented
second return value of that routine indicates containment in the alternating
group; it is \emph{not} a solvability flag \cite{sympyfields}. The implementation
uses the returned group's solvability property. This low-degree optimization
is not the degree range of the general algorithm.

\begin{proposition}\label{prop:base-solvable}
$H$ is solvable if and only if $G=\Gal(K/\Q)$ is solvable. Every prime divisor of
$|H|$ divides $M$.
\end{proposition}
\begin{proof}
Restriction identifies $H$ with $\Gal(K/K\cap C)$, a normal subgroup of $G$.
The quotient is $\Gal(K\cap C/\Q)$ and is abelian, because it is a quotient of
the abelian cyclotomic group. A group is solvable exactly when a normal
subgroup and its quotient are solvable. Also $|H|$ divides $|G|=D$, proving the
prime-divisor statement.
\end{proof}

Thus a nonsolvable $H$ is an actual impossibility result. Failure to construct
$H$ within a degree or time budget is not. When $H$ is solvable, compute
\begin{equation}\label{eq:series}
 H=H_0\triangleright H_1\triangleright\cdots\triangleright H_s=1,
 \qquad H_{i-1}/H_i\simeq C_{p_i}
\end{equation}
with $p_i$ prime. Only successive normality is needed; the subgroups need not
all be normal in $H$.

\section{Turning cyclic steps into radical generators}
\label{sec:resolvents}
Let
\[
 E_i=L^{H_i}.
\]
The fixed fields give an ascending tower
\[
 C=E_0\subset E_1\subset\cdots\subset E_s=L,
 \qquad [E_i:E_{i-1}]=p_i.
\]
Choose $\sigma\in H_{i-1}\setminus H_i$. Its restriction to $E_i$ generates the
cyclic quotient of order $p=p_i$. Its order as an automorphism of all of $L$
need not be $p$; what matters is that $\sigma^p$ is trivial on $E_i$. Put
$\epsilon=z^{M/p}\in C$.

\subsection{A deterministic nonzero Lagrange resolvent}
For $b\in E_i$, consider
\begin{equation}\label{eq:projector}
 \mathcal P(b)=\sum_{j=0}^{p-1}\epsilon^{-j}\sigma^j(b).
\end{equation}
Cyclic reindexing gives $\sigma(\mathcal P(b))=\epsilon\mathcal P(b)$.
Therefore, whenever $r=\mathcal P(b)\ne0$,
\begin{equation}\label{eq:kummer}
 r^p\in E_{i-1},\qquad E_i=E_{i-1}(r).
\end{equation}
The second assertion follows because $r$ is not fixed by $\sigma$ and the
relative degree is prime.

We need an actual finite search for $b$, not an instruction to ``choose a
generic element.'' For the power basis of $L$, form
\begin{equation}\label{eq:trace-candidates}
 b_k=\sum_{h\in H_i}h(t^k),\qquad 0\leq k<d.
\end{equation}
These elements span $E_i$ over $\Q$. Indeed, the trace onto $E_i$ is a
surjective $\Q$-linear map: on an element of $E_i$ it is multiplication by
$|H_i|$. Applying it to a basis of $L$ therefore gives a spanning set.

\begin{lemma}\label{lem:nonzero}
At least one of $\mathcal P(b_0),\ldots,\mathcal P(b_{d-1})$ is nonzero.
\end{lemma}
\begin{proof}
The distinct automorphisms $1,\sigma,\ldots,\sigma^{p-1}$ of $E_i$ are linearly
independent as maps with coefficients in $E_i$. For completeness, a minimal
nonzero dependence $\sum_j a_j\tau_j(x)=0$ for all $x$ can be evaluated at $ux$
and compared with $\tau_1(u)$ times the original dependence. Choosing $u$ on
which another $\tau_j$ differs from $\tau_1$ gives a shorter nonzero
dependence, a contradiction. Thus \eqref{eq:projector}, whose coefficients are
nonzero, is not the zero operator. It cannot vanish on a spanning set.
\end{proof}

The code scans the finite list in \eqref{eq:trace-candidates}, chooses the first
nonzero resolvent, and rescales it by a nonzero rational to normalize its
coefficient vector. It then checks exactly that $H_i$ fixes $r$, that
$\sigma(r)=\epsilon r$, and that every element of $H_{i-1}$ fixes $r^p$.
Rescaling affects the radicand by a rational $p$th power but changes neither
the eigenvalue nor the generated field.

\subsection{Coordinates by Fourier projection}
Merely finding $r_i$ is insufficient: one must express its radicand in earlier
generators and express the requested roots in the completed tower. There is a
particularly convenient coordinate formula. For $a\in E_i$ write
\[
 a=\sum_{j=0}^{p-1}c_jr^j,\qquad c_j\in E_{i-1}.
\]
Then
\begin{equation}\label{eq:coordinates}
 c_j=\frac1p\sum_{k=0}^{p-1}\sigma^k\!\left(\frac{a}{r^j}\right).
\end{equation}
To verify this, substitute the expansion of $a$. The coefficient of
$c_\ell r^{\ell-j}$ is the geometric sum
$p^{-1}\sum_k\epsilon^{k(\ell-j)}$, equal to one for $\ell=j$ and zero otherwise.
This also proves that each computed coefficient is in $E_{i-1}$.

Apply \eqref{eq:coordinates} recursively. At level zero, express an element of
$C$ in the rational basis $1,z,\ldots,z^{\varphi(M)-1}$ by exact linear algebra.
For every step this produces
\begin{equation}\label{eq:slp}
 r_i^{p_i}=A_i(z,r_1,\ldots,r_{i-1}),
 \qquad
 \alpha_j=B_j(z,r_1,\ldots,r_s).
\end{equation}
The coordinate expressions can be polynomials in the earlier generators with
rational coefficients, with exponents below the corresponding relative
degrees. The algorithm checks the reconstruction identity at each recursive
step. Memoization avoids recomputing the same coordinates, but no
simplification search is necessary for correctness.

\section{Coherent branches and selecting the intended root}
\label{sec:branches}
An abstract field generator is not born with a principal complex branch. It
would be incorrect to derive several unrelated identities, take the principal
root in each, and assume that all algebraic relations survive. For example,
principal square roots do not obey $\sqrt a\sqrt b=\sqrt{ab}$ for all complex
$a,b$. The remedy is a \emph{single shared tower}, not a collection of
independently simplified formulas.

\begin{theorem}[Coherent-branch root-set theorem]\label{thm:branches}
Suppose the exact relations in \eqref{eq:slp} describe the tower
$C=E_0\subset\cdots\subset E_s=L$, with each step having degree $p_i$.
Assign $z$ the principal value $(-1)^{2/M}$. In order, assign $r_i$ any one root
of
\[
 X^{p_i}-A_i(z,r_1,\ldots,r_{i-1}).
\]
Reuse that value everywhere $r_i$ occurs. Then the evaluated expressions
$B_1,\ldots,B_n$ are exactly the distinct roots of $f$. In particular, always
using the principal $p_i$th root is valid for the \emph{whole root set}.
\end{theorem}
\begin{proof}
The initial assignment extends to an embedding of $C$ into $\C$, because the
assigned value is a primitive $M$th root of unity. Assume an embedding of
$E_{i-1}$ has been defined. Since $E_i=E_{i-1}(r_i)$ has degree $p_i$, the
polynomial $X^{p_i}-A_i$ is its minimal polynomial. Applying the embedding to
its coefficients preserves irreducibility over the image field, and mapping
$r_i$ to any root of the image polynomial extends the embedding to $E_i$.
Inductively we obtain a field embedding $L\hookrightarrow\C$. It preserves all
polynomial identities and cannot identify two distinct roots. The $n$
distinct embedded roots of a degree-$n$ polynomial are the full root set.
\end{proof}

This result removes exponential enumeration of branch choices from the
construction. It does \emph{not} identify which output corresponds to the
user's root index. That is a separate comparison after the root set is known.
Nor does it justify a branch-unsafe transformation such as replacing the
principal radical of a product by a product of principal radicals.

\subsection{The Wolfram acceptance condition}
The Wolfram interface uses approximate distances only to order candidates.
It returns a candidate $e$ only after
\begin{lstlisting}
RootReduce[e - alpha] === 0
\end{lstlisting}
has succeeded, in addition to the radical-grammar check. A timeout or an
unevaluated comparison is not accepted. This prevents accidentally returning
a conjugate with the same minimal polynomial. No Python root index is assumed
to agree with Wolfram's root ordering.

\subsection{The Python selected-root comparison}
The Python entry point \code{solve\_root} uses zero-based \code{CRootOf}
ordering. It obtains the minimal factor of that particular root, constructs
its radical root set, and calls \code{Poly.same\_root}. The latter is documented
as a comparison between two roots of the supplied polynomial
\cite{sympyroots}. That precondition is important: it is supplied by the
algebraic construction, not inferred from numerical proximity.

In the tested SymPy implementation, this comparison uses a root-separation
bound and bounded-error evaluation rather than a fixed decimal tolerance.
The root-set certificate is an exact rational-arithmetic object; the
selected-index comparison additionally trusts SymPy's bounded-error
evaluation contract. We do not describe that numerical implementation as an
independently verified interval-arithmetic proof. Applications requiring a
smaller numerical trusted base can replace this final comparison with certified
complex-ball or rectangular root isolation, without changing the tower
construction.

\section{Positive certificates and their independent checker}
\label{sec:certificate}
The general construction returns more than an expression. Its certificate
contains a rational defining polynomial $P$ for the final field, the input
polynomial $f$, the integer $M$, coordinate vectors for $z$, the resolvents
$r_i$, and all roots of $f$, together with the coordinate expressions in
\eqref{eq:slp}. These data are serialized as integer strings and a small
expression syntax tree.

The checker is independent of the field-search history and of the Galois-group
computation. It performs the following checks in $F=\Q[t]/P$:
\begin{enumerate}
\item $P$ is irreducible, the polynomial being certified is nonconstant, and,
      when supplied, it agrees with the caller's expected polynomial up to
      monic normalization.
\item The proposed $z$ satisfies $\Phi_M(z)=0$. Each nonzero proposed $r_i$
      satisfies its stated power identity, using only earlier generators.
\item The monomials
\begin{equation}\label{eq:tower-basis}
 z^a r_1^{b_1}\cdots r_s^{b_s},\qquad
 0\leq a<\varphi(M),\quad 0\leq b_i<p_i,
\end{equation}
      number exactly $\deg P$ and have linearly independent rational coordinate
      vectors.
\item The output coordinates evaluate to the supplied distinct root witnesses,
      their number is $\deg f$, and each witness satisfies $f$ exactly.
\end{enumerate}

\begin{proposition}[Soundness of the positive checker]\label{prop:certificate}
If these checks succeed, evaluating the stated radical tower with coherent
principal branches produces the full set of roots of the certified polynomial.
\end{proposition}
\begin{proof}
Irreducibility of $P$ ensures that the ambient quotient is a field. The
cyclotomic relation gives a base field of degree $\varphi(M)$. Each power
relation bounds the next relative degree by $p_i$. The full independence of
\eqref{eq:tower-basis} forces equality at every step: a smaller degree at any
stage would make the final monomial family dependent. Thus the power relations
are genuine minimal polynomials for the successive generators. The branch
argument of Theorem~\ref{thm:branches} applies. The exact output and root
identities then give $\deg f$ distinct roots.
\end{proof}

The checker need not prove that $F$ is the smallest possible splitting field,
or even reconstruct its automorphisms. A larger valid radical field containing
all roots would be an equally sound positive certificate. It also does not
certify an external root index: that comparison remains separate.

The general constructor supports \code{verify=True}, and the Wolfram bridge
requests that check by default. Fast structural results are justified by their
recognized identities and ordinary degree-at-most-four solvers; they do not
currently carry this number-field certificate. A \status{NotSolvable} result
relies on the exact group calculation, not on a separately exported negative
certificate. These trust boundaries are deliberate and documented.

\section{Completeness, termination, and cost}
\label{sec:complete}
\begin{theorem}[Unbounded constructive algorithm]\label{thm:complete}
For an irreducible polynomial over $\Q$, the general procedure described above,
with resource bounds removed and correct terminating exact arithmetic
primitives, terminates. It either constructs radical expressions for all roots
or determines that none of its roots has a radical expression over $\Q$.
\end{theorem}
\begin{proof}
The splitting-field adjunction procedure terminates by Section
\ref{sec:closure}; each primitive-element change terminates by
Lemma~\ref{lem:primitive}. Cyclotomic adjunction is finite. Factoring the final
defining polynomial and constructing the automorphism table are finite exact
operations. Finite-group solvability testing and a composition series terminate.
If the group is nonsolvable, Theorem~\ref{thm:criterion} and
Proposition~\ref{prop:base-solvable} give the negative conclusion. Otherwise
there are finitely many prime steps; each has a nonzero resolvent in the finite
list of Lemma~\ref{lem:nonzero}. Coordinate recursion decreases its level at
every recursive call. The resulting tower gives the roots by
Theorem~\ref{thm:branches}.
\end{proof}

This theorem is not a claim that every program run with the default settings
will succeed. By default the Python backend limits a working field to degree
48 and requests a 120-second budget. These are engineering choices, not a
mathematical degree limit on radical-solvable inputs. A field of large degree
may appear even for a low-degree polynomial; a cyclotomic base change can make
it larger still. Exact coefficients may grow rapidly during primitive-element
changes, group construction, and resolvent coordinates.

The algorithm explicitly enumerates a group and constructs a full splitting
field. It is consequently not a polynomial-time implementation in the input
bit length. This should not be misrepresented as an inherent lower bound on
radical solving. Landau and Miller proved polynomial-time solvability testing
and a polynomial-time radical construction in straight-line-program form
\cite{landau-miller}. Implementing that substantially more sophisticated
algorithm is a different engineering project. Here the priority is an explicit,
auditably constructive reference algorithm with modest dependencies and useful
fast paths. Large expanded formulas can also destroy the advantage of a compact
straight-line representation, regardless of the underlying decision algorithm.

The finite-budget result \status{ResourceLimit} carries no solvability
conclusion. A native-only or fast-only miss is \status{NotFound}. An expression
that could not be matched to the requested root in time gives
\status{VerificationFailed} in the Wolfram interface. None of these should be
converted into \status{NotSolvable} by calling code.

\section{Software organization and use}
\label{sec:software}
\subsection{Files and dependencies}
The archive contains the article source and PDF, \code{SystematicRadicals.wl},
the Python backend, runnable examples, a regression suite, the actual Python
validation report, and saved positive certificates. The reference environment
was Python 3.13.5 with SymPy 1.14.0 on Linux. The dependency file pins SymPy to
that version because the implementation uses its exact algebraic-field
representations. Future versions should be revalidated rather than silently
assumed compatible. Python 3.10 or later is required by the source syntax.

No SageMath, PARI/GP, Mathematica installation, or network service is needed to
run the Python implementation. The Wolfram native layer needs no Python for
its fast paths; its general method invokes the bundled local backend.

\begin{lstlisting}
cd SystematicRadicals
python -m pip install -r requirements.txt
python examples/quickstart.py
python tests/run_validation.py
\end{lstlisting}

\subsection{Python API}
The main all-roots routine accepts an irreducible polynomial:
\begin{lstlisting}[language=Python]
from systematic_radicals import solve_polynomial, verify_certificate
import sympy as s
x = s.symbols("x")

answer = solve_polynomial(x**3 - 2, x,
    method="galois", max_field_degree=48,
    max_seconds=120, verify=True)
print(answer.assignments)
print(answer.outputs)
assert verify_certificate(answer.certificate, x**3 - 2)
\end{lstlisting}
Add the archive's \code{python} directory to \code{sys.path}, or use the example
scripts, which locate it automatically. \code{assignments} is the ordered
radical program; \code{outputs} contains expressions in its variables.
\code{expanded\_roots()} substitutes all assignments and can be expensive.
\code{diagnostics} records field degrees, group orders, and prime steps.

The default \code{method="auto"} tries identities before the general method.
\code{method="galois"} forces the general construction except for rational
roots, making it useful for tests.

The setting \code{method="fast"} never constructs a splitting field. To request the ideal unbounded construction use
\begin{lstlisting}[language=Python]
answer = solve_polynomial(f, x, method="galois",
    max_field_degree=None, max_seconds=None, verify=True)
\end{lstlisting}
This is an explicit opt-in to potentially extreme computation.

For a selected root, including an input polynomial with several irreducible
factors, use
\begin{lstlisting}[language=Python]
from systematic_radicals import solve_root
f = x**6 + x**4 - x**3 - x**2 - 1
positive_root = solve_root(f, 1, x)  # zero-based CRootOf index
\end{lstlisting}
Do not transfer indices from Wolfram Language. Exceptions distinguish
\code{InvalidInput}, \code{NotFound}, \code{NotSolvable},
\code{ResourceLimit}, and \code{CertificateError}. The Python library's time
budget is cooperative inside construction: an individual SymPy factorization
cannot be interrupted by a check that occurs only between operations. It does
not bound all preprocessing or selected-root evaluation. The command-line
worker described below provides the harder operational boundary.

\subsection{Wolfram Language API}
After placing the archive's directory at a convenient location:
\begin{lstlisting}
Get["/path/to/SystematicRadicals/SystematicRadicals.wl"];
alpha = Root[#^6 + #^4 - #^3 - #^2 - 1 &, 2];
r = Radicalize[alpha, Method -> "Native"];
VerifyRadical[alpha, r]
\end{lstlisting}
\code{Radicalize} returns an expression or \code{Failure}. \code{RadicalReport}
returns an association containing the expression, a pure-function minimal
polynomial, the method, and diagnostics. Its success flag is set only after
exact selected-root verification, except for an input already in the required
radical grammar, whose identity is immediate.

The principal options are:
\begin{center}\small
\begin{tabular}{@{}p{0.33\textwidth}p{0.18\textwidth}p{0.41\textwidth}@{}}
\toprule
Option & Default & Meaning\\
\midrule
\code{Method} & \code{Automatic} & Native methods, then Python fallback.\\
\code{"PythonExecutable"} & \code{Automatic} & Local interpreter path.\\
\code{"TimeLimit"} & $120$ & Requested overall computation budget.\\
\code{"VerificationTimeLimit"} & $30$ & Upper limit per exact candidate comparison.\\
\code{"MaxFieldDegree"} & $48$ & Cap on the backend's working field degree.\\
\code{"VerifyCertificate"} & \code{True} & Check a general positive certificate.\\
\bottomrule
\end{tabular}
\end{center}
\code{Method -> "Native"} forbids Python; \code{"Python"} skips the native
layer but allows Python fast paths; \code{"Galois"} forces the general backend
unless the input is already a radical expression.
An explicit Python executable path is useful on Windows or with a virtual
environment. The process is invoked as an argument list with \code{RunProcess},
not through shell string interpolation \cite{runprocess}.

For an unbounded attempt set all three relevant bounds to \code{Infinity}:
\begin{lstlisting}
Radicalize[alpha, Method -> "Galois",
  "MaxFieldDegree" -> Infinity,
  "TimeLimit" -> Infinity,
  "VerificationTimeLimit" -> Infinity]
\end{lstlisting}
Startup, shutdown, operating-system scheduling, and external-library calls mean
that the Wolfram option is not a real-time deadline guarantee. Backend failure,
protocol failure, mathematical nonsolvability, and inconclusive verification
have different failure tags. The supplied \code{.wlt} file can be run with
\code{TestReport}; those Wolfram tests have not been executed in this work.

\subsection{Data-only JSON and the worker boundary}
The backend can be called directly:
\begin{lstlisting}
python python/systematic_radicals.py --json
\end{lstlisting}
It reads one JSON object from standard input. Coefficients are listed from
highest to lowest degree, each as a numerator/denominator pair of decimal
integer strings. For example:
\begin{lstlisting}
{"coefficients": [["1","1"],["0","1"],["0","1"],["-2","1"]],
 "method":"galois", "max_field_degree":48,
 "max_seconds":120, "verify":true, "certificate":true}
\end{lstlisting}
The response contains a status, diagnostics, ordered assignments, and a root
list. Expressions use only the tags \code{Q}, \code{Ref}, \code{Add},
\code{Mul}, and \code{Pow}. References are zero-based indices into earlier
assignments. The Wolfram decoder validates the tags and integer strings and
never calls \code{ToExpression} on subprocess data. Invalid or forward
references are rejected.

A separate spawned worker performs construction and certificate checking. Its
parent waits for the requested time and terminates the worker on expiry.
Process startup and shutdown add overhead; the timeout applies to the worker
wait, not to every millisecond of the command's lifetime. A successful
command exits with code zero; other statuses use exit code two. This design
keeps an expensive exact factorization from indefinitely blocking a bounded
command-line request, subject to ordinary operating-system process control.

\section{Executed validation and remaining limitations}
\label{sec:validation}
All 36 Python tests passed in the recorded run. The full run took about
81.38 seconds. The test log and machine-readable report are included in
\code{validation}; the table below gives selected test durations, including the
checks performed by each test, rather than claiming isolated algorithmic
benchmarks.

\begin{center}\small
\begin{tabular}{@{}p{0.40\textwidth}p{0.38\textwidth}r@{}}
\toprule
Input or test & Method/check & Seconds\\
\midrule
Original sextic, all roots & Reciprocal reduction, root matching & 0.51\\
Original quintic, all roots & Dickson identity, root matching & 0.03\\
$(x-3)^7-2$ & Power composition, root matching & 5.93\\
$x^6+8x^4+x^3+16x^2+8$ & Weighted reciprocal, root matching & 1.14\\
$x^2-2$ & Forced tower, independent certificate & 0.03\\
$x^3-2$ & Forced tower, independent certificate & 0.43\\
$x^3-3x+1$ & Forced tower, independent certificate & 0.30\\
$x^4-10x^2+1$ & Forced tower, independent certificate & 0.12\\
$x^4-2$ & Forced nonabelian tower, certificate & 22.17\\
$x^3-x-1$ & Forced nonabelian tower, certificate & 37.62\\
$x^5-x-1$ & Exact nonsolvability & 0.01\\
\bottomrule
\end{tabular}
\end{center}

The forced general examples are significant even though several have elementary
closed forms: they exercise the actual field-and-group path instead of allowing
a cubic or quartic formula to conceal an unimplemented fallback. For
$x^3-x-1$, the splitting field has degree six, the cyclotomic compositum has
degree twelve, and the stabilizer composition orders are $6,3,1$. For
$x^4-2$, the corresponding orders are $8,4,2,1$. The latter is a nonabelian
2-group case; the former requires both a quadratic and a cubic cyclic step.
Saved certificates permit repeating the independent positive checks without
redoing the field search.

The suite also checks malformed inputs, nonrational coefficients, a root
selected from a reducible polynomial, rational and complex radical grammar,
field-degree limits, a hard worker timeout, altered radicands, duplicate root
witnesses, and a certificate presented for the wrong polynomial. The tests of
all-root matching rely on already established polynomial membership, as
required by \code{same\_root}. They are not a substitute for the certificate
proof.

For the nonsolvable control there is a small independent explanation. Modulo
two,
\[
 x^5-x-1=(x^2+x+1)(x^3+x^2+1),
\]
and the factors are distinct and irreducible. Modulo three the polynomial is
irreducible. Thus it is irreducible over $\Q$ and its transitive degree-five
Galois group contains an element of cycle type $(2)(3)$. Cubing that element
gives a transposition. A transitive group of prime degree is primitive; a
primitive permutation group containing a transposition is the full symmetric
group. One way to see the latter assertion is to take the graph whose edges
are the conjugate transpositions: its connected components form a block
system, so the graph is connected, and the edge transpositions generate the
full symmetric group. Therefore the group is $S_5$, which is not solvable.
The program's recorded group order is $120$, with derived subgroup of order
$60$.

There are substantial performance limitations. The general construction is
coefficient-height sensitive and can be slow beyond small working fields.
An exploratory forced construction for a cyclic quintic was stopped while an
internal exact operation was still running; it is not counted as a successful
test or included in the timing table. This illustrates why cooperative library
checks and externally enforced worker bounds are distinct. No claim is made
that the general path has been stress-tested across all solvable transitive
groups, all degrees, or all branch-degenerate-looking expressions. The native
Wolfram interface has only been statically inspected, not kernel-tested.

The code also does not optimize radical depth, the number of radicals, total
expression size, or the size of the cyclotomic base beyond the squarefree
choice. It does not implement arbitrary polynomial decomposition, arbitrary
Tschirnhaus transformations, or a real-radical-only variant. These are absent
optimizations or different specifications, not hidden assumptions in the
completeness theorem for complex radicals over $\Q$.

\section{Conclusions and a development path}
A systematic solver needs more than a powerful simplifier. It needs a precise
algebraic domain, a genuine solvability test, a construction that reaches every
solvable case in principle, and a branch-aware test for the selected root.
The supplied implementation separates those responsibilities.

For the motivating examples, exact structural identities produce small answers
without expensive Galois arithmetic. When these identities fail, the general
fallback actually constructs the normal field, its automorphisms, and its
cyclic prime steps. The Lagrange projectors have a finite nonzero search, the
coordinate recursion is explicit, and the resulting positive certificate is
checkable without trusting the construction's group search. Coherent principal
branches are enough to obtain the complete root set; exact root identification
then handles the original \code{Root} object.

The most valuable future improvements are performance-oriented: smaller
primitive elements, better intermediate-field choices, a compact program
representation preserved across the Wolfram bridge, and a more sophisticated
Galois implementation avoiding full regular representations. A separately
checkable negative certificate and an independently certified root-index
comparison would further reduce the trusted computational base. None of those
is silently claimed by version 1.0.0. What is supplied is a working exact
reference construction, useful native fast paths, reproducible positive
certificates, and explicit boundaries on the guarantees of a bounded run.

\appendix
\section{Algorithm-to-code map}
\begin{center}\small
\begin{longtable}{@{}p{0.34\textwidth}p{0.58\textwidth}@{}}
\toprule
Component & Python implementation\\
\midrule
\endhead
Input normalization & \code{\_poly}, \code{solve\_root}\\
Structural identities & \code{\_fast}, \code{dickson}, \code{\_rational\_roots\_of\_power}\\
Abstract field arithmetic & \code{NumberField}\\
Primitive-element adjunction & \code{NumberField.adjoin}\\
Splitting field & \code{\_normal\_closure}\\
Cyclotomic base and automorphisms & First part of \code{\_tower}\\
Composition series and resolvents & Middle of \code{\_tower}\\
Recursive Fourier coordinates & Nested \code{express} function in \code{\_tower}\\
Straight-line radical output & \code{RadicalSolution}, \code{encode\_expr}\\
Positive checker & \code{verify\_certificate}\\
Selected root & \code{solve\_root}\\
JSON protocol and bounded worker & \code{\_request}, \code{\_worker}, \code{main}\\
\bottomrule
\end{longtable}
\end{center}
The Wolfram layer exposes two solving functions, \code{Radicalize} and
\code{RadicalReport}, and two validation functions, \code{RadicalExpressionQ}
and \code{VerifyRadical}.
The Python backend is the executed reference implementation for the general
algorithm; the Wolfram file is an interface and an independent implementation
of the inexpensive structural layer.

\section{Input provenance and reproducibility}
The supplied archive was inspected as data and was not executed. Its SHA-256
hash is
\begin{lstlisting}
381e80382e8736aa00e18d8fd48c6f391b9ba0b15e80ff7621e7589147174930
\end{lstlisting}
The extracted \code{FindExtension.nb} has SHA-256
\begin{lstlisting}
d70951fdb88128556ee7ea762652413b6f8b5e970c4437c2f3293d9b07eac007
\end{lstlisting}
The original notebook is not redistributed in this package. The new code is
provided under the MIT license. The included article is provided under the
license stated in the archive. Build the PDF from this directory with two runs
of \code{pdflatex SystematicRadicals.tex}, or with \code{latexmk -pdf}.

\begin{thebibliography}{9}
\bibitem{question}
Vladimir Reshetnikov.
\emph{Why is ToRadicals[] not able to handle all cases? Is there a workaround?}
Mathematica Stack Exchange, question 34011, 2013; consulted 8 September 2026.
\url{https://mathematica.stackexchange.com/questions/34011/why-is-toradicals-not-able-to-handle-all-cases-is-there-a-workaround}.

\bibitem{milne}
J. S. Milne.
\emph{Fields and Galois Theory}, version 5.10, 2022.
In particular, the cyclic-extension and solvability results in Section 5.
\url{https://www.jmilne.org/math/CourseNotes/FT.pdf}.

\bibitem{landau-miller}
Susan Landau and Gary Lee Miller.
Solvability by radicals is in polynomial time.
\emph{Journal of Computer and System Sciences} \textbf{30}(2) (1985), 179--208.
\url{https://www.cs.cmu.edu/~glmiller/Publications/Papers/LaMi85.pdf}.

\bibitem{toradicals}
Wolfram Research.
\emph{ToRadicals---Wolfram Language Documentation}.
Consulted 8 September 2026.
\url{https://reference.wolfram.com/language/ref/ToRadicals.html}.

\bibitem{runprocess}
Wolfram Research.
\emph{RunProcess---Wolfram Language Documentation}.
Consulted 8 September 2026.
\url{https://reference.wolfram.com/language/ref/RunProcess.html}.

\bibitem{sympyfields}
SymPy Development Team.
\emph{Number Fields}, SymPy 1.14.0 documentation.
Consulted 8 September 2026.
\url{https://docs.sympy.org/latest/modules/polys/numberfields.html}.

\bibitem{sympygroups}
SymPy Development Team.
\emph{Permutation Groups}, SymPy 1.14.0 documentation.
Consulted 8 September 2026.
\url{https://docs.sympy.org/latest/modules/combinatorics/perm_groups.html}.

\bibitem{sympyroots}
SymPy Development Team.
\emph{Polynomial Manipulation Module Reference}, in particular
\code{Poly.same\_root}, SymPy 1.14.0 documentation.
Consulted 8 September 2026.
\url{https://docs.sympy.org/latest/modules/polys/reference.html}.
\end{thebibliography}
\end{document}
